\documentclass[11pt,a4paper]{article}
\usepackage[margin=1in]{geometry}
\usepackage[T1]{fontenc}
\usepackage[utf8]{inputenc}
\usepackage{amsmath,amssymb}
\usepackage{booktabs}
\usepackage{tabularx}
\usepackage{hyperref}
\usepackage{setspace}
\usepackage{enumitem}
\usepackage{array}
\usepackage{caption}
\usepackage{xurl}
\setstretch{1.05}
\hypersetup{colorlinks=true,linkcolor=black,citecolor=black,urlcolor=blue}
\setlength{\parskip}{0.45em}
\setlength{\parindent}{1.2em}
\setlist[itemize]{leftmargin=1.4em,itemsep=0.15em}
\setlist[enumerate]{leftmargin=1.4em,itemsep=0.15em}
\newcommand{\sha}[1]{{\footnotesize\ttfamily\path{#1}}}

\title{Block Order under a Fixed English--Tagalog Token Budget:\\
\large A One-Seed Schedule-Topology Study with nanochat}
\author{Paul Pajo\\
De La Salle -- College of Saint Benilde\\
\texttt{paulamerigo.pajojr@benilde.edu.ph}}
\date{25 August 2026}

\begin{document}
\maketitle

\begin{abstract}
Holding English--Tagalog source-content token quotas, tokenizer, parent type, phase-two budget, and optimizer-reset policy fixed, does phase-two \emph{schedule topology}---the order and granularity of language blocks---change Tagalog and English full-split validation bits-per-byte (BPB)? This question was preregistered as AsPredicted \#307969 after P5 Gate~W and before any P6 parent, child, validation BPB, or restricted-test BPB existed. P6-M is a post-P5, prospectively preregistered one-seed schedule-topology mechanism study. It does not amend AsPredicted \#306780, \#306935, \#307342, \#307591, or \#307836. It is not a P5-style recurrence-count study. One parent-initialization seed ($s=4$) trained a fresh Tagalog parent on WikiText-TL-39 with Karpathy's nanochat trainer pinned at commit \texttt{92d63d4}, the carry-forward P4 tokenizer, $T=2048$, and NVIDIA CUDA A40. After P0-T PASS, d20 was frozen as immutable C0. Six matched d20 continuations compared C1 extra Tagalog, C2 pure English, and four mixed topologies (M-fine, M-coarse, M-blocked, M-rand) that each use exactly $9{,}633{,}792$ Tagalog and $9{,}633{,}792$ English source-content tokens for $D_{\mathrm{phase2}}=19{,}267{,}584$ tokens ($N=294$). Primary estimands are $\Delta\mathrm{TL}(\tau)=\mathrm{TL}(M\textrm{-}\tau)-\mathrm{TL}(\textrm{M-fine})$ and $\Delta\mathrm{EN}(\tau)=\mathrm{EN}(M\textrm{-}\tau)-\mathrm{EN}(\textrm{M-fine})$ for $\tau\in\{\textrm{coarse},\textrm{blocked},\textrm{rand}\}$ at $\delta=0.01$. At Gate~X, M-coarse was better than M-fine on Tagalog and worse on English; M-blocked was worse on both; M-rand was within $\delta$ on both. One authorized M-fine-only secondary test is reported separately and excluded from topology classification. One seed; no confidence interval, mean, $p$-value, or population claim.
\end{abstract}

\noindent\textbf{Keywords:} schedule topology; block order; token budget; Tagalog; English; bits-per-byte; nanochat; preregistration; post-P5; one seed.

\noindent\textbf{arXiv categories:} cs.CL, cs.LG.

\noindent\textbf{Registration and deposits (P6-M):} AsPredicted \#307969 (\url{https://aspredicted.org/bk6m9d.pdf}); study record \url{https://github.com/pageman/nanochat-filipino/tree/main/docs/papers/p6-m-schedule-topology}; GitHub trees \texttt{scripts/p6/}, \texttt{docs/run-cards/p6/}, \texttt{manifests/p6/}, \texttt{docs/hub/p6-m-schedule-topology/} under \url{https://github.com/pageman/nanochat-filipino}; provisional Hub ID \texttt{pageman/nanochat-filipino-p6-m-schedule-topology} (upload deferred). ResearchBox and AsCollected deposit IDs are deferred at Gate~W closeout. Run ID: \texttt{p6-20260824T155226Z-769f807a}. Manuscript version: v1.0.

\section{Introduction}
\label{sec:intro}

Pajo~\cite{pajo2026p11} measured equal-exposure Tagalog BPB across nanochat depths on WikiText-TL-39 (AsPredicted \#306780). Pajo~\cite{pajo2026p2} asked the forward continual-pretraining question from an English parent (AsPredicted \#306935). Pajo~\cite{pajo2026p3} asked the reverse from a fresh Tagalog parent (AsPredicted \#307342). Pajo~\cite{pajo2026p4} locked source-content token share at $q_{\mathrm{TL}}=0.50$ after a fresh Tagalog parent (AsPredicted \#307591). Pajo~\cite{pajo2026p5} counted how often that frozen mix recurred on three unused initializations (AsPredicted \#307836).

P6-M was designed after those released P1.1--P5 findings were known. It is therefore a prospectively preregistered \textbf{post-P5} schedule-topology mechanism study~\cite{aspredicted307969}. It is not an amendment, correction, rescue, replication, or independent confirmation of P4 or P5. It does not treat recurrence counts as primary. It holds the P4 token-share bundle fixed and varies only the filed block schedules.

The confirmatory question, locked before any P6 confirmatory BPB existed, is:

\begin{quote}
After a newly trained seed-4 Tagalog parent passes P0-T, freeze d20 as C0 and continue it for one shared phase-two token budget on C1, C2, and four mixed topologies that share identical within-language streams and exact per-language quotas. Relative to M-fine, how do M-coarse, M-blocked, and M-rand classify on Tagalog and English validation BPB at $\delta=0.01$?
\end{quote}

This paper contributes: (i)~one eligible seed-4 Tagalog parent frozen as C0; (ii)~six sibling children under a fixed token budget; (iii)~a sealed twelve-cell validation matrix before one M-fine-only secondary test; (iv)~one Gate~X release whose primary result is the six primary $\Delta$TL/$\Delta$EN classifications versus M-fine.

\section{Related work}
\label{sec:related}

Cruz and Cheng constructed WikiText-TL-39~\cite{cruz2019eval}. Sequential interference is classical~\cite{mccloskey1989catastrophic,french1999cf}. LLM continual-pretraining work often continues an English or multilingual base~\cite{luo2023forgetting,ramasesh2021forgetscale,shi2024continual}. P6-M stays inside WikiText-family text, the nanochat depth dial, and BPB rather than CORE or chat~\cite{karpathy2026nanochat}. Equal token budgets follow the family so stream identity is not confounded with compute~\cite{kaplan2020scaling,hoffmann2022chinchilla,muennighoff2023dataconstrained}. WikiText established clean Wikipedia LM evaluation in English~\cite{merity2017pointer}. BPB remains comparable across vocabularies~\cite{shannon1951prediction,brown2020gpt3}. Official P6 BPB is mean token NLL divided by $\ln 2$ times mean UTF-8 bytes per token, full split, $T=2048$, BOS-best-fit. In-loop trainer BPB is not confirmatory. Preregistration reduces undisclosed flexibility~\cite{nosek2018prereg,wagenmakers2012confirmatory}. Curriculum learning motivates asking whether presentation order matters when exposure totals are locked~\cite{bengio2009curriculum}; surveys of continual LLM training likewise treat schedule and interference as open mechanism questions~\cite{shi2024continual}. P6-M answers the block-order question only inside this one-seed BPB apparatus. WikiText-TL-39 is a Cruz--Cheng resource lineage; this study is not an extension of Cheng catastrophic-forgetting methods and does not assess downstream Filipino tasks~\cite{cruz2020degradation}.

\section{Methods}
\label{sec:methods}

\subsection{Registration and authority}
AsPredicted \#307969 governs~\cite{aspredicted307969}.
\begin{itemize}[leftmargin=1.6em,itemsep=0.35em]
\item PDF: \url{https://aspredicted.org/bk6m9d.pdf}
\item PDF SHA-256:\\\sha{769f807a00264a996b02c38b83ee2cf6c23c4981e36fb477dfc1959fa918a1e7}
\item Prefiling addendum SHA-256:\\\sha{df49664809ada69d23dcd2c799e75b30f0fdb9afd7aee12b071ac24ef2f81082}
\item Unmodified gate-plan SHA-256:\\\sha{d8a63608608c59d2c4d9882e5346625462056c331094942a8a01d496697a1c79}
\item Topology-manifest SHA-256:\\\sha{d1e7d5af7247a572e319ee003b5f4e3da5d1fb1592e5ed9ff6b22eeec15ea606}
\end{itemize}
Authority order: filed PDF $\gg$ SHA-bound addendum $\gg$ unmodified gate plan $\gg$ \texttt{LOCK.json}, manifests, and deterministic receipts $\gg$ this manuscript (v1.0). The study does not amend \#306780, \#306935, \#307342, \#307591, or \#307836. ResearchBox and Hub uploads were deferred at Gate~W; local archive and paper hashes are recorded in the run-card closeout.

\subsection{Sources, tokenizer, and hosts}
Tagalog train/val/test documents are the frozen WikiText-TL-39 P1.1-eligible splits. English train/val/test documents are the frozen WikiText-103 P2-eligible splits. The tokenizer is the carry-forward P4 pair: \texttt{tokenizer.pkl} \sha{04436b854e0841025a3dd2b46baaeeea07a7ccc252e9f99a19171306f00bc5a8} and \texttt{token\_bytes.pt} \sha{a5dbc1c88f6292696108263072d77115718cc2d8357f7ad4859adfa517cc2132}. P1.1--P5 weights were not loaded as parents. Official GPU gates used NVIDIA CUDA A40. Official evaluation used \texttt{scripts/p6/evaluate\_bpb.py} with formula-equivalent hash \sha{a22a566b71ac9b768543e052e78576f6b079d5c88066f3ec2f8bfef600e17481}.

\subsection{Parent, P0-T, and C0 freeze}
Parent-initialization seed is exactly $s=4$. After P0-T PASS, d20 was copied to an immutable C0 directory with \texttt{additional\_train\_tokens}=0. C0 SHA-256: \sha{7fbd24de792aa4ee27d841866db2114e0fb45b1fdcaa4edcc9b24582220123c9}. Smoke and d8 checkpoints are not C0.

\subsection{Children and fixed topology treatments}
C1, C2, M-fine, M-coarse, M-blocked, and M-rand each continued the same frozen C0 for $N=294$ steps, $B=65{,}536$, $D_{\mathrm{phase2}}=19{,}267{,}584$ tokens, serial R$\to$S$\to$T1--T4, \texttt{load\_optimizer=False}, child peak learning rate $0.3\times$ the parent peak, warmup 14, terminal checkpoint only. Mixed arms share identical within-language streams and stored no-wrap schedules:
\begin{itemize}
\item \textbf{M-fine:} EN-first alternating blocks of $2{,}048$ tokens ($4{,}704$ blocks per language).
\item \textbf{M-coarse:} EN-first alternating blocks of $1{,}204{,}224$ tokens (eight blocks per language).
\item \textbf{M-blocked:} all Tagalog then all English (one block per language).
\item \textbf{M-rand:} one precomputed shuffle of $4{,}704$ EN and $4{,}704$ TL blocks of length $2{,}048$, shuffled once with Python \texttt{random.Random(42)} and hash-pinned before filing.
\end{itemize}
Each mixed arm uses exactly $9{,}633{,}792$ Tagalog and $9{,}633{,}792$ English source-content tokens. Mix-manifest SHA-256: \sha{f203c615266bc8c33c358c1de397715791cae33536a9743c8a6bf8cd543cb107}.

\subsection{Lockbox, analysis, and tests}
Validation BPB was lockboxed until one Gate~X. Gate~U sealed twelve child cells (six conditions $\times$ two languages) with \texttt{test\_access}=0. Policy~A authorized exactly one M-fine-only secondary restricted-test event after U (Gate~V). C1, C2, M-coarse, M-blocked, and M-rand were not tested. Primary contrasts use sealed validation cells; tests are secondary and excluded from topology classification. Equality outside $(-\delta,+\delta)$ at $\delta=0.01$ defines directional classes. No composite score, mean, confidence interval, $p$-value, across-seed average, or P5-style recurrence count.

\section{Results}
\label{sec:results}

P0-T status is \textbf{PASS} for seed 4. All six children completed with reload-verified terminal checkpoints. Table~\ref{tab:cells} reports the sealed twelve validation cells and descriptive C0 English. Table~\ref{tab:primary} reports the six primary contrasts versus M-fine. Table~\ref{tab:contextual} reports contextual $R_{\mathrm{TL}}$ and $A_{\mathrm{EN}}$ for the four mixed topologies.

\begin{table}[ht]
\centering
\caption{Full-split validation BPB after terminal d20 checkpoints. Six-decimal reporting as filed. C0 English is descriptive and excluded from topology classification.}
\label{tab:cells}
\begin{tabular}{lcc}
\toprule
Arm & Tagalog val BPB & English val BPB \\
\midrule
C0 (frozen parent) & --- & $2.683814$ \\
C1 extra Tagalog & $0.564548$ & $2.892551$ \\
C2 pure English & $2.219915$ & $1.269724$ \\
M-fine & $1.256292$ & $1.440084$ \\
M-coarse & $1.214675$ & $1.487389$ \\
M-blocked & $1.458898$ & $1.529111$ \\
M-rand & $1.257379$ & $1.441159$ \\
\bottomrule
\end{tabular}
\end{table}

\begin{table}[ht]
\centering
\caption{Primary contrasts versus M-fine at $\delta=0.01$. $\Delta=\mathrm{BPB}(M\textrm{-}\tau)-\mathrm{BPB}(\textrm{M-fine})$. Lower BPB is better.}
\label{tab:primary}
\begin{tabular}{lcccc}
\toprule
$\tau$ & $\Delta$TL & TL class & $\Delta$EN & EN class \\
\midrule
M-coarse & $-0.041617$ & better than M-fine by $\delta$ & $0.047305$ & worse than M-fine by $\delta$ \\
M-blocked & $0.202606$ & worse than M-fine by $\delta$ & $0.089027$ & worse than M-fine by $\delta$ \\
M-rand & $0.001087$ & within $\delta$ & $0.001075$ & within $\delta$ \\
\bottomrule
\end{tabular}
\end{table}

\begin{table}[ht]
\centering
\caption{Contextual contrasts (descriptive secondary). $R_{\mathrm{TL}}=\mathrm{TL}(M\textrm{-}\tau)-\mathrm{TL}(\mathrm{C2})$; $A_{\mathrm{EN}}=\mathrm{EN}(M\textrm{-}\tau)-\mathrm{EN}(\mathrm{C1})$.}
\label{tab:contextual}
\begin{tabular}{lcc}
\toprule
$\tau$ & $R_{\mathrm{TL}}$ & $A_{\mathrm{EN}}$ \\
\midrule
M-fine & $-0.963623$ & $-1.452467$ \\
M-coarse & $-1.005239$ & $-1.405162$ \\
M-blocked & $-0.761016$ & $-1.363439$ \\
M-rand & $-0.962536$ & $-1.451392$ \\
\bottomrule
\end{tabular}
\end{table}

Under this preregistered fixed-budget schedule comparison, M-rand was within $\delta$ of M-fine on both languages; M-blocked was worse than M-fine by more than $\delta$ on both languages; M-coarse was better than M-fine on Tagalog and worse on English, each by more than $\delta$. This is a one-seed topology classification, not a confidence interval, not a population effect, and not a P5 recurrence count. P6-M does not amend \#307591 or \#307836.

One authorized M-fine-only secondary test event reported English test BPB $1.450423$ and Tagalog test BPB $1.260475$. Tests do not alter sealed contrasts and are not co-primary.

\section{Discussion}
\label{sec:disc}

With token share held fixed, block order was not inert in this apparatus: coarse interleaving and complete blocking moved validation BPB relative to fine EN-first alternation, while a once-shuffled random block sequence stayed within $\delta$ of M-fine. That pattern is a mechanism observation under a frozen P4-style quota bundle, not a claim that M-fine is optimal, not a general forgetting-mitigation result, and not a law across seeds. Contextual $R_{\mathrm{TL}}$ and $A_{\mathrm{EN}}$ place each mixed arm relative to C2 and C1; they are not a second topology success rule. WikiText-TL-39 remains a Cruz--Cheng resource lineage; P6-M does not extend Cheng catastrophic-forgetting protocols and does not claim downstream Filipino-task or all-Philippine-language benefit.

\section{Limitations}
\label{sec:limit}

One seed; one tokenizer; WikiText-family text only; no chat, SFT, CORE, FilBench, or larger-model transfer. Confirmatory GPU work used a rented NVIDIA A40. Raw test JSONL remains restricted. Hub weight release, if any, must ship C0+C1+C2+topology children together with tokenizer and meta; never overwrite P4/P5 Hub IDs. Document-revisit (P7-M), optimizer-state (P8-M), replay/protection, and tokenizer/transfer studies remain separate filings.

\section{Deviations}
\label{sec:dev}

No protocol stop. No break-glass. No additional confirmatory validation or test eval after Gate~X. Gate~S attempt 1 failed as a \emph{preflight} block when the C2 English packed stream was missing on the pod path ($checkpoint\_sha256=\texttt{null}$); after syncing the filed stream, one clean authorized C2 launch completed. M-rand wrote a valid terminal model checkpoint, then hit ENOSPC during optimizer save; the terminal model was hash- and reload-verified and accepted without retraining (\texttt{technical\_accept=true}). Gate~0 had pinned a P5-rename unblind script hash that reads \texttt{c3\_*} cells and emits recurrence counts; Gate~U wrote topology-named cells, so Gate~X executed a topology-contrast script under filed-analysis authority. The script-hash mismatch is documented in the unblinding event; the filed primary remains $\Delta$TL/$\Delta$EN versus M-fine. Frozen English train JSONL contains intra-split duplicate rows ($28{,}472$ rows / $28{,}232$ unique); they were not dropped. Cross-split overlap was 0.

\section*{Availability}
\begin{itemize}[leftmargin=1.6em,itemsep=0.35em]
\item AsPredicted \#307969: \url{https://aspredicted.org/bk6m9d.pdf}
\item Study record (GitHub):\\
\url{https://github.com/pageman/nanochat-filipino/tree/main/docs/papers/p6-m-schedule-topology}
\item Code and audit trees (GitHub):\\
\url{https://github.com/pageman/nanochat-filipino}\\
P6-only: \texttt{scripts/p6/}, \texttt{docs/papers/p6-m-schedule-topology/},
\texttt{docs/run-cards/p6/}, \texttt{manifests/p6/},
\texttt{docs/hub/p6-m-schedule-topology/}
\item Weights (Hugging Face Hub): provisional ID\\
\url{https://huggingface.co/pageman/nanochat-filipino-p6-m-schedule-topology}\\
(upload deferred at Gate~W; never write onto P4/P5 Hub IDs).
\item ResearchBox / AsCollected: deferred at Gate~W closeout.
\end{itemize}
Held-out test text is not redistributed. Host credentials are not published. GitHub holds protocol, receipts, sealed/released JSON, and paper; Hugging Face is reserved for optional weight bundles after independent remote re-hash.

\section*{Ethics}
Public Wikipedia-derived corpora. No human-subjects experiment. Any future ResearchBox remains passcode-protected for peer review until explicitly made public.

\section*{Acknowledgements}
Thanks to Manus 1.6 and Cursor.com(Auto) for the drafting, formatting, and solutioning of the paper. WikiText-TL-39 is due to Cruz and Cheng~\cite{cruz2019eval}. WikiText-103 is due to Merity et al.~\cite{merity2017pointer}. nanochat is due to Karpathy~\cite{karpathy2026nanochat}. Errors remain the author's.

\section*{Funding}
No dedicated grant. GPU time (Runpod A40) was paid by the author.

\appendix
\section{Selected hashes}
\label{app:hashes}
\begin{itemize}
\item Pin / nanochat: \sha{92d63d4e8bb4df75c3b71618f31ddde2378b2bcd}
\item AsPredicted PDF: \sha{769f807a00264a996b02c38b83ee2cf6c23c4981e36fb477dfc1959fa918a1e7}
\item Gate plan: \sha{d8a63608608c59d2c4d9882e5346625462056c331094942a8a01d496697a1c79}
\item Addendum: \sha{df49664809ada69d23dcd2c799e75b30f0fdb9afd7aee12b071ac24ef2f81082}
\item Topology manifest: \sha{d1e7d5af7247a572e319ee003b5f4e3da5d1fb1592e5ed9ff6b22eeec15ea606}
\item Mix manifest: \sha{f203c615266bc8c33c358c1de397715791cae33536a9743c8a6bf8cd543cb107}
\item Evaluator (P6 path): \sha{a22a566b71ac9b768543e052e78576f6b079d5c88066f3ec2f8bfef600e17481}
\item C0: \sha{7fbd24de792aa4ee27d841866db2114e0fb45b1fdcaa4edcc9b24582220123c9}
\item C1 / C2: \sha{6223c116779ce3128c1e4cae0f2e03744b3068169ad7bb88f75a5146671a99bd} / \sha{04c06c195e513c7b30752637b80591c29e740879d91c35c323fb75fafca8747d}
\item M-fine / M-coarse / M-blocked / M-rand: \sha{a0139607a8fdf2772d4b8e722b449b6a8db04056a8dc38cd177708af8f15eeab} / \sha{86588ee9f72ea4a85a821c2c882a363d55d3e50ccb325bdabaa0706e1e911dfe} / \sha{5ec45216eb0a09b5bc04f04f4622d85f7d8f7ff8861fdea6de5487f7c18fa526} / \sha{38580cd501c0d87a1a502397697f7c7e427134eb7cb64bdce2ba6be1a036f108}
\item U seal: \sha{23f9a94483532f5b7d51f4a27f77436cafd440f30c11fac5e0ea2dfa0aa095ae}
\item Released contrasts: \sha{d376643af19969f2ee127afcc4696de3a4821c8188fa6e7003d190d5f4031911}
\end{itemize}

\section{GitHub versus Hugging Face}
\label{app:split}
GitHub \texttt{pageman/nanochat-filipino} is the study record: scripts, paper, lock, ledgers, run-card receipts, sealed/released JSON, and Hub \emph{documentation}. Hugging Face \texttt{pageman/nanochat-filipino-p6-m-schedule-topology} is the optional weight deposit (deferred): C0+C1+C2+topology children plus tokenizer and evaluation JSON. Never a single topology child alone. Never write onto P1.1/P2/P3/P4/P5 Hub IDs. Raw test JSONL, passcodes, SSH keys, optimizer states, and HOST operator cards belong in neither public tree.

\begin{thebibliography}{99}

\bibitem{pajo2026p11}
P.~Pajo.
\newblock Equal-exposure depth and held-out Tagalog bits-per-byte on WikiText-TL-39.
\newblock AsPredicted \#306780, 2026.

\bibitem{pajo2026p2}
P.~Pajo.
\newblock Held-out English bits-per-byte after matched-budget Tagalog continuation of a WikiText-103 nanochat parent.
\newblock AsPredicted \#306935, 2026.

\bibitem{pajo2026p3}
P.~Pajo.
\newblock Tagalog retention and English acquisition under equal-budget nanochat continual pretraining.
\newblock AsPredicted \#307342, 2026.

\bibitem{pajo2026p4}
P.~Pajo.
\newblock Token-share-locked English--Tagalog mixtures after a fresh Tagalog parent.
\newblock AsPredicted \#307591, 2026.

\bibitem{pajo2026p5}
P.~Pajo.
\newblock How often a frozen mix recurs across unused initializations: a closed three-seed count, not a population estimate.
\newblock AsPredicted \#307836, 2026.

\bibitem{aspredicted307969}
AsPredicted \#307969.
\newblock P6-M: schedule topology under a frozen English--Tagalog token-share bundle.
\newblock \url{https://aspredicted.org/bk6m9d.pdf}, 2026.

\bibitem{cruz2019eval}
J.~C.~B.~Cruz and C.~Cheng.
\newblock Evaluating Filipino linguistic resources.
\newblock 2019.

\bibitem{cruz2020degradation}
J.~C.~B.~Cruz and C.~Cheng.
\newblock Establishing baselines for text classification in low-resource languages.
\newblock \emph{arXiv:2005.02068}, 2020.
\newblock Discussion-only robustness analogy for P6-M; not the BPB schedule-topology estimand.

\bibitem{karpathy2026nanochat}
A.~Karpathy.
\newblock nanochat.
\newblock \url{https://github.com/karpathy/nanochat}, 2026.

\bibitem{merity2017pointer}
S.~Merity, C.~Xiong, J.~Bradbury, and R.~Socher.
\newblock Pointer sentinel mixture models.
\newblock \emph{ICLR}, 2017.

\bibitem{shi2024continual}
H.~Shi et al.
\newblock Continual learning of large language models: A comprehensive survey.
\newblock \emph{arXiv:2404.16789}, 2024.

\bibitem{luo2023forgetting}
Y.~Luo et al.
\newblock An empirical study of catastrophic forgetting in large language models during continual fine-tuning.
\newblock \emph{arXiv:2308.08747}, 2023.

\bibitem{ramasesh2021forgetscale}
V.~V.~Ramasesh, A.~Lewkowycz, and E.~Dyer.
\newblock Effect of scale on catastrophic forgetting in neural networks.
\newblock \emph{ICLR}, 2022.

\bibitem{mccloskey1989catastrophic}
M.~McCloskey and N.~J.~Cohen.
\newblock Catastrophic interference in connectionist networks: The sequential learning problem.
\newblock \emph{Psychology of Learning and Motivation}, 24:109--165, 1989.

\bibitem{french1999cf}
R.~M.~French.
\newblock Catastrophic forgetting in connectionist networks.
\newblock \emph{Trends in Cognitive Sciences}, 3(4):128--135, 1999.

\bibitem{kaplan2020scaling}
J.~Kaplan et al.
\newblock Scaling laws for neural language models.
\newblock \emph{arXiv:2001.08361}, 2020.

\bibitem{hoffmann2022chinchilla}
J.~Hoffmann et al.
\newblock Training compute-optimal large language models.
\newblock \emph{NeurIPS}, 2022.

\bibitem{muennighoff2023dataconstrained}
N.~Muennighoff et al.
\newblock Scaling data-constrained language models.
\newblock \emph{NeurIPS}, 2023.

\bibitem{shannon1951prediction}
C.~E.~Shannon.
\newblock Prediction and entropy of printed English.
\newblock \emph{Bell System Technical Journal}, 30(1):50--64, 1951.

\bibitem{brown2020gpt3}
T.~Brown et al.
\newblock Language models are few-shot learners.
\newblock \emph{NeurIPS}, 2020.

\bibitem{nosek2018prereg}
B.~A.~Nosek et al.
\newblock The preregistration revolution.
\newblock \emph{PNAS}, 115(11):2600--2606, 2018.

\bibitem{wagenmakers2012confirmatory}
E.-J.~Wagenmakers et al.
\newblock An agenda for purely confirmatory research.
\newblock \emph{Perspectives on Psychological Science}, 7(6):632--638, 2012.

\bibitem{bengio2009curriculum}
Y.~Bengio, J.~Louradour, R.~Collobert, and J.~Weston.
\newblock Curriculum learning.
\newblock \emph{ICML}, 2009.

\end{thebibliography}

\end{document}
